%%%%%%%%%%%%%%%%%%%%%%%%%%%%%%%%%%%%%%%%%%%%%%%%%%%%%%%%%%%%%%%%%%%%%%%%%%%%%%%%
\chapter{Covariant Statistical Mechanics in the SHP Formalism}
%%%%%%%%%%%%%%%%%%%%%%%%%%%%%%%%%%%%%%%%%%%%%%%%%%%%%%%%%%%%%%%%%%%%%%%%%%%%%%%%

\section{Introduction}

The preceding chapters established the Hamiltonian formulation of the
Stueckelberg--Horwitz--Piron theory together with its gauge-covariant
generalization. The dynamics of an individual particle is completely
determined by Hamilton's equations in the extended phase space.

Physical systems, however, rarely consist of a single isolated particle.
Macroscopic phenomena arise from large ensembles whose microscopic
states cannot be specified individually. A statistical description is
therefore essential.

The purpose of this chapter is to formulate statistical mechanics
directly on the SHP phase space. The resulting theory differs from
ordinary relativistic statistical mechanics in several important
respects.

First, evolution is parametrized by the invariant world time
\(\tau\), rather than by the coordinate time \(x^0\).

Second, the invariant mass is not constrained to remain constant.
Consequently, statistical ensembles possess a distribution not only in
space-time and momentum but also in off-shell states.

Third, if internal non-Abelian charges are present, the statistical
state must also include the corresponding color degrees of freedom.

The natural microscopic state of the system is therefore specified by

\[
(x^\mu,p_\mu,Q^a),
\]

while the statistical state is described by the probability density

\[
\boxed{
\rho(x^\mu,p_\mu,Q^a,\tau).
}
\]

The remainder of this chapter develops the corresponding Liouville
theorem, kinetic equations, equilibrium theory and entropy directly from
the Hamiltonian structure established in the previous chapters.



%%%%%%%%%%%%%%%%%%%%%%%%%%%%%%%%%%%%%%%%%%%%%%%%%%%%%%%%%%%%%%%%%%%%%%%%%%%%%%%%
\section{Probability Distribution on SHP Phase Space}

Consider an ensemble of identical systems.

Each microscopic state corresponds to one point of the extended phase
space

\[
\Gamma
=
T^*M
\times
{\cal O},
\]

where \({\cal O}\) denotes the internal color manifold.

The statistical state of the ensemble is described by the distribution
function

\[
\boxed{
\rho
=
\rho(x^\mu,p_\mu,Q^a,\tau).
}
\]

The probability that the system occupies an infinitesimal phase-space
volume

\[
d\Gamma
\]

is

\[
dP
=
\rho\,d\Gamma.
\]

Normalization requires

\[
\boxed{
\int_\Gamma
\rho
\,d\Gamma
=
1.
}
\]

Unlike ordinary relativistic kinetic theory, the SHP distribution is
defined on the complete off-shell phase space.

Consequently, particles possessing different invariant masses coexist
within the same statistical ensemble.


%%%%%%%%%%%%%%%%%%%%%%%%%%%%%%%%%%%%%%%%%%%%%%%%%%%%%%%%%%%%%%%%%%%%%%%%%%%%%%%%
\section{Statistical Distribution of the Invariant Mass}
%%%%%%%%%%%%%%%%%%%%%%%%%%%%%%%%%%%%%%%%%%%%%%%%%%%%%%%%%%%%%%%%%%%%%%%%%%%%%%%%

In ordinary relativistic statistical mechanics the mass shell

\[
p^\mu p_\mu = m^2
\]

is imposed as a kinematic constraint before the statistical theory is
constructed. Consequently, the invariant mass does not constitute an
independent statistical variable.

The SHP formalism adopts a fundamentally different viewpoint. Since the
Hamiltonian dynamics does not constrain particles to remain on a fixed
mass shell, the invariant mass becomes a dynamical quantity whose value
may differ from one member of the ensemble to another.

It is therefore natural to regard the invariant mass as an additional
statistical variable.

%%%%%%%%%%%%%%%%%%%%%%%%%%%%%%%%%%%%%%%%%%%%%%%%%%%%%%%%%%%%%%%%%%%%%%%%%%%%%%%%

\subsection{Mass Distribution Function}

Let

\[
\Pi_\mu
=
p_\mu-gA_\mu^aQ^a
\]

denote the gauge-covariant momentum.

The corresponding invariant is

\[
\mu^2
=
\Pi^\mu\Pi_\mu .
\]

The probability density of finding a particle with invariant mass
\(\mu\) is defined by

\[
\boxed{
P(\mu^2,\tau)
=
\int_\Gamma
\rho(x,p,Q,\tau)\,
\delta\!\left(
\mu^2-\Pi^\alpha\Pi_\alpha
\right)
d\Gamma.
}
\]

This distribution satisfies the normalization condition

\[
\boxed{
\int
P(\mu^2,\tau)
\,d\mu^2
=
1.
}
\]

Thus the statistical ensemble possesses a complete spectrum of
off-shell states.

%%%%%%%%%%%%%%%%%%%%%%%%%%%%%%%%%%%%%%%%%%%%%%%%%%%%%%%%%%%%%%%%%%%%%%%%%%%%%%%%

\subsection{Statistical Moments}

The first statistical moment defines the mean invariant mass,

\[
\boxed{
\langle\mu^2\rangle
=
\int
\mu^2
P(\mu^2,\tau)
\,d\mu^2.
}
\]

Higher moments characterize the shape of the distribution.

For example,

the variance is

\[
\boxed{
\sigma_\mu^2
=
\langle(\mu^2)^2\rangle
-
\langle\mu^2\rangle^2.
}
\]

Unlike ordinary relativistic mechanics,

these quantities generally evolve with the world time.

%%%%%%%%%%%%%%%%%%%%%%%%%%%%%%%%%%%%%%%%%%%%%%%%%%%%%%%%%%%%%%%%%%%%%%%%%%%%%%%%

\subsection{Evolution Equation}

Differentiating the definition of

\[
P(\mu^2,\tau)
\]

with respect to

\[
\tau,
\]

and using the Liouville equation,

\[
\frac{\partial\rho}{\partial\tau}
+
\{\rho,K\}
=
0,
\]

one finds

\[
\frac{\partial P}{\partial\tau}
=
-
\frac{\partial}{\partial\mu^2}
J_\mu,
\]

where

\[
J_\mu
\]

is the probability current in invariant-mass space.

Thus probability is conserved not only in ordinary phase space but also
along the invariant-mass direction.

%%%%%%%%%%%%%%%%%%%%%%%%%%%%%%%%%%%%%%%%%%%%%%%%%%%%%%%%%%%%%%%%%%%%%%%%%%%%%%%%

\subsection{Stationary Mass Spectrum}

If

\[
\frac{\partial P}{\partial\tau}=0,
\]

then

\[
J_\mu
=
\text{constant}.
\]

For a closed ensemble with vanishing current at the boundaries,

\[
J_\mu=0,
\]

and the mass spectrum becomes stationary.

Stationarity therefore corresponds to equilibrium in invariant-mass
space.

%%%%%%%%%%%%%%%%%%%%%%%%%%%%%%%%%%%%%%%%%%%%%%%%%%%%%%%%%%%%%%%%%%%%%%%%%%%%%%%%

\subsection{Physical Interpretation}

The existence of the distribution

\[
P(\mu^2,\tau)
\]

is a direct consequence of off-shell dynamics.

Different members of the statistical ensemble occupy different values of
the invariant mass, and the Hamiltonian evolution continuously
redistributes probability among these values.

The mass shell of ordinary relativistic mechanics is therefore replaced
by a dynamical mass spectrum.

The invariant mass becomes a statistical observable rather than a fixed
kinematic parameter.



